
\documentclass[aspectratio=169]{beamer}

\mode<presentation>
\usetheme{Boadilla}
\definecolor{NTHU1}{RGB}{127,16,132}
\definecolor{NTHU2}{RGB}{228,210,231}
\definecolor{blue}{RGB}{30,90,205}
\definecolor{red}{RGB}{213,94,0}
\definecolor{green}{RGB}{0,128,0}
\definecolor{charcoalgray}{RGB}{108,104,104}
\setbeamercolor{title}{fg=NTHU1}
\setbeamercolor{frametitle}{fg=NTHU1}
\setbeamercolor{block title}{bg=NTHU1, fg=white}
\setbeamercolor{block body}{bg=NTHU2}
\setbeamercolor{structure}{fg=NTHU1}
\setbeamercolor{item projected}{fg=white}
\setbeamercolor{item}{fg=NTHU1}
\setbeamercolor{subitem}{fg=NTHU1}
\setbeamercolor{section in toc}{fg=NTHU1}
\setbeamercolor{description item}{fg=NTHU1}
\setbeamercolor{caption name}{fg=NTHU1}
\setbeamercolor{button}{bg=NTHU1, fg=white}
\setbeamercolor{caption name}{fg=NTHU1}
\usepackage{graphics}
\usepackage{tikz}
\usepackage{amsmath}
\usepackage{bbm}
\usetikzlibrary{decorations.pathreplacing}
\usepackage{geometry}
\usepackage{booktabs}
\usepackage{bookmark}
\usepackage{multirow, makecell}
\usepackage{float}
\usepackage{fancyvrb}
\usepackage{caption}
\captionsetup[figure]{singlelinecheck = false, format= hang, justification=centering}
\captionsetup[subfigure]{singlelinecheck = false, format= hang, justification=raggedright}
\usepackage{subcaption}
\usepackage{adjustbox}
\usepackage{hyperref}
\usepackage{threeparttable}
\usepackage{appendixnumberbeamer} 
\usepackage[T1]{fontenc}
\usepackage[default]{lato} 
\usepackage{smartdiagram}
\smartdiagramset{
  back arrow disabled = true,
  module minimum width=1cm,
  module x sep = 3,
  uniform arrow color=true,
}
\usepackage{xcolor}
\usepackage{colortbl}
	\definecolor{light-gray}{gray}{0.99}

\newenvironment{wideitemize}{\itemize\addtolength{\itemsep}{10pt}}{\enditemize}
\newenvironment{wideenumerate}{\enumerate\addtolength{\itemsep}{10pt}}{\endenumerate}
\newenvironment{widedescription}{\description\addtolength{\itemsep}{10pt}}{\enddescription}
\hypersetup{
colorlinks=true,
linkcolor=NTHU1,
filecolor=green, 
urlcolor=blue,
}
\beamertemplatenavigationsymbolsempty
\setbeamercolor{author in head/foot}{bg=white, fg=NTHU1}
\setbeamercolor{title in head/foot}{bg=white, fg=NTHU1}
\setbeamercolor{date in head/foot}{bg=white, fg=NTHU1}
\setbeamercolor{section in head/foot}{bg=white, fg=NTHU1}
\setbeamercolor{headline}{bg=white}
\setbeamertemplate{footline}{
    \leavevmode%
    \hbox{%
        \begin{beamercolorbox}[wd=.333333\paperwidth,ht=2.25ex,dp=1ex,center]{date in head/foot}%
            \usebeamerfont{date in head/foot}%\insertdate
        \end{beamercolorbox}%
        \begin{beamercolorbox}[wd=.333333\paperwidth,ht=2.25ex,dp=1ex,center]{title in head/foot}%
            \usebeamerfont{title in head/foot}\insertshorttitle
        \end{beamercolorbox}%
        \begin{beamercolorbox}[wd=.333333\paperwidth,ht=2.25ex,dp=1ex,right]{date in head/foot}%
            \usebeamerfont{date in head/foot}\insertframenumber{}\hspace*{2em}
        \end{beamercolorbox}}%
        \vskip0pt%
    }
    %% May need to script the introduction!!!!!!
%\setbeamercolor{page number in head/foot}{fg=NTHU1}
\setbeamertemplate{section in toc}[sections numbered]
\setbeamertemplate{subsection in toc}{\leavevmode\leftskip=3em\rlap{\hskip-1.75em\inserttocsectionnumber.\inserttocsubsectionnumber}
\inserttocsubsection\par}
\setbeamerfont{subsection in toc}{size=\footnotesize}
%\setbeamertemplate{headline}{%
  %\begin{beamercolorbox}[ht=5.5ex]{section in head/foot}
    %\vskip2pt\insertnavigation{0.33\paperwidth}\vskip2pt
  %\end{beamercolorbox}%
%}
\newenvironment{transitionframe}{\setbeamercolor{background canvas}{bg=white}\setbeamertemplate{footline}{} \begin{frame}[noframenumbering]}{\end{frame}}


\makeatletter
\let\@@magyar@captionfix\relax
\makeatother


\title[]{Public Economics and Development} % Change this regularly
\subtitle[]{Week 1: Course Introduction}
\author[Lee]{Seung-hun Lee }
\institute[]{Taipei School of Economics\\ National Tsing Hua University}

\date[]{February 24th, 2026}

\begin{document}
\begin{transitionframe}
\titlepage
\end{transitionframe}


%%% Color slides for section headers: Use for colloquium version (The ones bewteen \iffals and \fi)

\section{Course logistics}
\begin{frame}
\frametitle{Course logistics} 
Instructor: Seung-hun Lee 

\begin{itemize}\setlength\itemsep{3pt}

\item Office location: Room C06 in this building
\item Office hours: Tuesdays 2:00-4:00PM \href{https://www.calendly.com/seunghunlee0918}{(Sign up here)}
\item Ph.D. in Economics from Columbia University (2024)
\end{itemize}\medskip
Class time and location
\begin{itemize}
\item Time: Tuesdays 9AM - 12PM
\item Location: Room A09
\end{itemize}\medskip
Coming prepared to the lecture
\begin{itemize}
\item Lecture notes are uploaded a day in advance 
\item Slides summarize textbooks and papers -- no need to purchase them separately
\item Keep up with the readings: more on that later
\end{itemize}\medskip
\end{frame}

\begin{frame}
\frametitle{Course goal: Understand public economics in developing countries} 
Study government's role in economic life: Taxation, public goods, and bureaucracy
\begin{itemize}\setlength\itemsep{3pt}
\item Classical theories on optimal taxation,  redistribution, and public sector employment
\item Income taxes, antipoverty programs, principal-agent problems etc.
\end{itemize}\medskip

 Understand challenges faced by developing countries
\begin{itemize}\setlength\itemsep{3pt}
\item Lagging behind on fiscal, legal, and personnel capacity of the state  
  \item Thus, what's optimal elsewhere do not work on developing countries
\end{itemize}\medskip 

Goals and requirements
\begin{itemize}\setlength\itemsep{3pt}
\item Be able to critically engage with scholarly work
\item Conduct independent, original research
\item Familiarity with statistical inference, micro/macroeconomics is necessary!
\end{itemize}
\end{frame}


\begin{frame}
\frametitle{Course goal: Understand public economics in developing countries} 
Main texts
\begin{itemize}\setlength\itemsep{3pt}
\item Atkinson, Stiglitz (2015): Lectures on Public Economics
\item Besley, Persson (2009): Pillars of Prosperity
\item Hindriks, Myles (2013): Intermediate Public Economics
\end{itemize}\medskip

 Recommended (and often helpful texts)
\begin{itemize}\setlength\itemsep{3pt}
\item Public economics: Gruber (2022, undergrad), Roland (2014, undergrad), Salanié (2011, graduate)
\item Econometrics: Angrist, Pischke (2009, graduate), Cameron, Trivedi (2005, graduate), Stock, Watson (2020, undergraduate)
\end{itemize}\medskip 

Readings
\begin{itemize}\setlength\itemsep{3pt}
\item 4-6 papers from leading economic scholarly journals
\item Read up on the starred papers in the reading list!
\end{itemize}
\end{frame}


\begin{frame}
\frametitle{Course topics (Refer to syllabus for full reading list)} 

Weeks 1-2: Introduction and econometric methods
\begin{itemize}\setlength\itemsep{3pt}
\item Introduction to public economics
\item Empirical analysis tools: OLS, IV, Panel regression, Natural experiments
\end{itemize}\medskip

Week 3-5: Classical theories on taxation and public economics
\begin{itemize}\setlength\itemsep{3pt}
\item How drastic are problems with inequality and welfare across countries and time?
\item How are optimal taxation and public expenditure achieved?
\item Why are these difficult to achieve in developing countries?
\end{itemize}\medskip

 Week 6-8: Fiscal capacity in developing countries
\begin{itemize}\setlength\itemsep{3pt}
\item What explains different tax collection patterns in developed and developing countries?
\item How does large informal economy affect taxation? 
\item How are developing countries adopting their tax structure to increase fiscal capacities?
\end{itemize}

\end{frame}

\begin{frame}
\frametitle{Course topics (Refer to syllabus for full reading list)} 
Week 9-11: Anti-poverty programs and public provision
\begin{itemize}
\item How should anti-poverty programs be designed?
\item Form of anti-poverty programs: Cash v. in-kind? Conditional or not?
\item How should public expenditures be targeted?
\end{itemize}\medskip
Week 12-14: Organizational capacity of the state
\begin{itemize}
\item  How to optimally allocate authority across bureaucrats and managers?
\item How can public sector recruit capable and motivated individuals?
\item How do public sector incentivize and monitor its workers?
\end{itemize}\medskip

\end{frame}




\begin{frame}
\frametitle{Tasks and grading scale} 

 Two referee reports (20\% of the grades)
\begin{itemize}\setlength\itemsep{3pt}
\item Aim: Critically engage with scholarly work (Guides: \href{https://papers.ssrn.com/sol3/papers.cfm?abstract_id=2547191}{here} and \href{https://scholar.harvard.edu/files/apassalacqua/files/refguidelines_tepe.pdf}{here}) 
  \item 200-word summary + 1500-word critique of not-yet published academic papers
\item Assess strength and weaknesses, while separating critical matters from minor issues 
\item For papers to referee, refer to the syllabus
\end{itemize}\medskip
 Three Problem sets (20\% of the grades) - Only the top two scores counts towards grade
\begin{itemize}\setlength\itemsep{3pt}
\item Solve questions testing theoretical problem-solving skills and empirical analysis skills
\item Theory: Solve mathematical model and identify intuitions in plain english
\item Empirics: Replicate and extend parts of an empirical scholarly paper
\end{itemize}
\end{frame}

\begin{frame}
\frametitle{Tasks and grading scale} 

 Research Proposal and Virtual Paper (30\% of the grades)
\begin{itemize}\setlength\itemsep{3pt}
\item Develop a viable plan for a paper and build-up towards a full paper
\item 2-page research proposal outlining research questions, expected contribution, background, potential data sources and empirical strategy (10\% of grade)
\item Virtual paper: Extend the proposal towards a paper with original analysis, building up to a full paper (20\% of grade)\\ \smallskip
$\Rightarrow$ \textbf{Unfinished but original and exciting project $>>$ complete, yet unintersting project}
\end{itemize}\medskip
 Final exam (30\% of the grades) 
\begin{itemize}\setlength\itemsep{3pt}
\item Cumulative exam at the end of the semester
\item Tests understanding of theoretical frameworks and connect it to empirical findings
\item Exact mode (take-home vs in-person) TBD
\end{itemize}
\end{frame}
\begin{frame}
\frametitle{More on the virtual paper} 

Format of the virtual paper: Mimics many scholarly articles
\begin{enumerate}\setlength\itemsep{3pt}
\item \underline{Introduction}: State your research question and why it matters (academically / policy)
\item \underline{Literature review}: State the gap in the relevant literature, and how you fill that gap
\item \underline{Background}: State the relevant information on the context that you are studying
\item \underline{Data}: Provide explanation on data source and why it works in your favor
\item \underline{Empirical strategy}: Explain research methods, whether it allows for causal inference, and potential limitations in that method
\item \underline{Preliminary results}: Highlight main findings (causational relationships)
\item \underline{Mechanism and falsification tests}: State what explains your main results, and what plans you have to address potential factors that may nullify your findings
\end{enumerate}
\end{frame}

\begin{frame}
\frametitle{Important deadlines} 
\begin{itemize}\setlength\itemsep{10pt}
\item March 31st: Referee report 1 due
\item April 14th: Problem set 1 due
\item April 28th: Research proposal due
\item May 12th: Problem set 2 due
\item May 19th: Refere report 2 due
\item June 2nd: Virtual paper due
\item June 9th: Problem set 3 due 
\item Week of June 9th: Final exam (mode to be declared later)
\end{itemize}
\end{frame}


\begin{frame}
\frametitle{Changes in schedule around March}
No class on March 24th and March 31st (Conferences in Europe and Thailand) 
\begin{itemize}\setlength\itemsep{3pt}
\item Two make up classes will be scheduled
\item One on March 30th, likely evening
\item One on the week of April 6th (perhaps friday)
\item Open to other schedules
\end{itemize}
\end{frame}




 
\begin{frame}
\frametitle{Some ground rules} 
Active discussions and questions!\par \medskip
 Usage of AI: Strictly for assistive purposes
\begin{itemize}\setlength\itemsep{3pt}
\item Allowed for proofreading, coding support
\item Do not generate complete assignments using AI 
\item All writings, especially virtual paper, must be original
\item To prove originality, regular discussion of your ideas will suffice
\end{itemize} \medskip
Rules regarding collaboration
\begin{itemize}\setlength\itemsep{3pt}
\item You are not allowed to collaborate on referee report!
\item Others: Collaboration is allowed, but under following conditions 
\item Each student must write up their own assignments
\item Students must clearly identify their collaborators
\end{itemize}
\end{frame}



\section{What is Public Economics?}
\begin{transitionframe}
\begin{center}
  \begin{huge}
    \textcolor{NTHU1}{%
What is Public Economics?}
  \end{huge}
\end{center}
\end{transitionframe}

\begin{frame}
\frametitle{Public economics study role of the state in our economic lives} 
 State collects taxes from its citizens, and provides many goods and services
\begin{itemize}\setlength\itemsep{3pt}
\item  Income tax, sales tax, social security contributions, excise tax... 
\item Education services, infrastructure, public safety, national defense, social safety net...
\item Significant number of population works at the public sector
\end{itemize} \medskip
\textit{Positive} analysis: Explanation of economic phenomena following government actions 
\begin{itemize}\setlength\itemsep{3pt}
\item How do taxes affect (or distort) labor supply, savings, or other economic behaviors? 
\item How are the effects distributed across different types of individuals?
\end{itemize} \medskip
\textit{Normative} analysis: Making judgements on how the government should act 
\begin{itemize}\setlength\itemsep{3pt}
\item How progressive should the tax systme be?
\item To whom should we provide more public services?
\end{itemize} \medskip
In this course, we focus on three pillars of public economics: Taxation, public goods, and personnel
\end{frame}




\begin{frame}
\frametitle{Why do taxes matter} 
Sources of revenue for governments to finance their key services
\begin{itemize}\setlength\itemsep{3pt}
\item Finance essential public goods and services (defense, infrastructure, education, health)
\item Reduce inequality through redistribution
\item May come at a cost: Distort labor incentives $\to$ Fall in efficiency  
\end{itemize} \medskip

Development of tax revenues and unresolved questions
\begin{itemize}\setlength\itemsep{3pt}
\item Tax share of GDP: 10\% in 1910 $\to$ 30-50\% today (in developed countries)
\item What is the optimal design of taxation?
\begin{itemize}\setlength\itemsep{3pt}
\item How progressive should taxes be?
\item Which taxes minimize distortions while raising needed revenue?
\end{itemize} 
\item Many developing countries lag behind: 
\begin{itemize}\setlength\itemsep{3pt}
\item What explains low tax collection capacity? 
\item Do optimal tax prescriptions from developed countries apply here?
\end{itemize} 
\end{itemize} \medskip
\end{frame}

\begin{frame}
\frametitle{Rise of the state in the developed countries} 
\begin{figure}
\centering
\includegraphics[keepaspectratio, width=0.8\textwidth]{fig1.png}
\end{figure}
\end{frame}

\begin{frame}
\frametitle{Why are taxes enforced differently?} 
Differences in observed policy choices across place and time 
\begin{itemize}\setlength\itemsep{3pt}
\item Significant progressive income taxes and sizeable public sector in developed countries
\item Many developing countries rely on regressive indirect taxes (consumption, tariffs, etc.)
\item Limited public services and lack of state presence in many fragile states
\end{itemize} \medskip

These differences originate from differences in underlying \textit{state capacity}
\begin{itemize}\setlength\itemsep{3pt}
\item Ability of states to implement policies, collect taxes, and provide public goods
\item Not all states can collect taxes digitally or have comprehensive information systems
\item Not all states are fully aware of citizens' income and can target policies accordingly
\item Not all states have sizeable public sectors capable of delivering policies and services
\end{itemize} \medskip

\end{frame}

\begin{frame}
\frametitle{Growth of fiscal (information) capacity in developed countries} 
\begin{figure}
\centering
\includegraphics[keepaspectratio, width=0.8\textwidth]{fig2.png}
\end{figure}
\end{frame}

\begin{frame}
\frametitle{.. that is not evenly shared across countries (Okunogbe, Tourek 2024)} 
\begin{figure}
\centering
\includegraphics[keepaspectratio, width=0.7\textwidth]{ot2024.png}
\caption*{Index of tax tech: Use of 3rd party info, tax ID, among other technology}
\end{figure}
\end{frame}

\begin{frame}
\frametitle{Thus, drastic difference in tax intakes (Besley, Persson 2014)} 
\begin{figure}
\centering
\includegraphics[keepaspectratio, width=0.75\textwidth]{fig3.png}
\end{figure}
\end{frame}

\begin{frame}
\frametitle{Why do public goods matter} 
Effective provision of public goods determines quality of life
\begin{itemize}\setlength\itemsep{3pt}
\item Support markets: Law and order that enforce contracts and protect property rights
\item Benefits to society beyond individual benefits: health, education, infrastructure etc.
\item Underprovided by private markets: They often under-internalize social benefits
\end{itemize} \medskip
Public goods provision is central to reduction in poverty and inequality
\begin{itemize}\setlength\itemsep{3pt}
\item The rich have means to seek private alternatives, or move to `better' locations
\item The poor frequently do not
\item Public spending serves to minimize discrepancy across the poor and the rich
\item Huge differences across how much money is spent on such programs across countries
\end{itemize} \medskip
\end{frame}



\begin{frame}
\frametitle{Expansion of public services over time in Europe} 
\begin{figure}
\centering
\includegraphics[keepaspectratio, width=0.75\textwidth]{fig4.png}
\end{figure}
\end{frame}

\begin{frame}
\frametitle{But significantly different even within developed countries} 
\begin{figure}
  \includegraphics[keepaspectratio, width=\textwidth]{oecd.png}
\caption{Public social expenditure as share of GDP}
\end{figure}
\end{frame}

\begin{frame}
\frametitle{Various public services in developing countries} 
\begin{figure}
\centering
\includegraphics[keepaspectratio, width=\textwidth]{fig5.png}
\end{figure}
\end{frame}

\begin{frame}
\frametitle{...still lag behind developed countries} 
\begin{figure}
\centering
\includegraphics[keepaspectratio, width=\textwidth]{fig6.png}
\end{figure}
\end{frame}

\begin{frame}
\frametitle{Unsettled questions in public goods / anti-poverty programs?} 
To whom do we target these programs? How much?
\begin{itemize}\setlength\itemsep{3pt}
\item Do you target specific populations? Or do you provide it universally?
\item Do you provide support in-kind or cash?
\item Do you attach conditions to income support or not?
\item How do you increase take-up for welfare programs? (surprisingly low in many places)
\end{itemize} \medskip
How do you improve welfare while minimizing side-effects?
\begin{itemize}\setlength\itemsep{3pt}
\item Labor supply potentially discouraged for both recipient and the rich
\begin{itemize}\setlength\itemsep{3pt}
\item Recipients: More earned income usually implies less support $\to$ work effort $\Downarrow$
\item Rich: More tax collected $\to$ less income from working $\to$ work effort $\Downarrow$
\end{itemize}
  \item How do you make welfare programs politically acceptable?
\item How do you design a program that also considers fiscal sustainability?
\end{itemize} \medskip
\end{frame}

\begin{frame}
\frametitle{Why public workers are important public economics} 
 Think of them as bureaucrats working for public institutions
\begin{itemize}\setlength\itemsep{3pt}
\item They are the frontline public service providers
\item They serve crucial roles for implementing various policies
\item They are sizeable: 35\% of all formal employment (World Bank data)
\end{itemize} \medskip
The quality of these workers matter, but also differ across contexts
\begin{itemize}\setlength\itemsep{3pt}
\item The human capital of these workers determine the quality of service provided
\item There are also large variation in the quality of public workers
\begin{itemize}\setlength\itemsep{3pt}
\item Tend to be `clustered' with economic development
\item Problems range from absenteeism, corruption, and lack of capabilities
\end{itemize}
\end{itemize} \medskip
\end{frame}

\begin{frame}
\frametitle{Public sector takes up huge share of employment globally} 
\begin{figure}
  \includegraphics[keepaspectratio, width=0.8\textwidth]{fig7.png}
\caption{Share of public sector employment in total employment, 2019 (Our World in Data)}
\end{figure}
\end{frame}

\begin{frame}
\frametitle{Development and bureaucratic quality co-move (Besley et al 2022)} 
\begin{figure}
  \includegraphics[keepaspectratio, width=0.65\textwidth]{fig8.png}
\caption*{Bureaucracy score: Avg of bureaucratic recruitment and impartial administrative practices}
\end{figure}
\end{frame}

\begin{frame}
\frametitle{Finding and retaining motivated and dedicated workers} 
 How do you persuade capable individuals to public sector?
\begin{itemize}\setlength\itemsep{3pt}
\item Individuals with high abilities AND prosocial motives are ideal targets
\item Public sectors in many contexts (not all!) provide lower and inflexible incentives
\item Thus, smaller applicant pool and less qualified Individuals (Dal Bó et al 2013)
\item Higher pay may attract individuals who apply `just for the money' (Dessserano 2015)
\end{itemize} \medskip
How do you incentivize and monitor existing public workers?
\begin{itemize}\setlength\itemsep{3pt}
\item Public workers' interests may not always align well with the public
\item Many places tried monitoring and incentive mechanisms, with mixed results
\item Incentives may increase effort, but also more abuse of discretions (Balan et al 2022)
\item Consistent monitoring is costly and potentially discourages efforts
\end{itemize} \medskip
\end{frame}


\begin{frame}
\frametitle{Quality of bureaucrat recruits vary across organizations} 
\begin{figure}
  \includegraphics[keepaspectratio, width=\textwidth]{indonesia.png}
\caption*{(Source: World Bank Group (2024), \textit{Innovating Bureaucracy for a More Capable Government})}
\end{figure}
\end{frame}


\begin{frame}
\frametitle{Task allocation not by merit: Possible disincentives} 
\begin{figure}
  \includegraphics[keepaspectratio, width=\textwidth]{WBG.png}
\caption*{(Source: World Bank Group (2024), \textit{Innovating Bureaucracy for a More Capable Government})}
\end{figure}
\end{frame}


\begin{frame}
\frametitle{We will often use mathematical models in our discussions} 

Models are a simiplification of reality that help us think about complex policy questions
\begin{itemize}\setlength\itemsep{3pt}
\item \textbf{Derive predictions}: What happens to people's behavior when we change tax rates? 
\item \textbf{Guide empirics}: What should we measure? What comparisons are informative?
\item \textbf{Welfare analysis}: Does the policy make society better off/equitable?
\end{itemize} \medskip
Most models in this course share common elements:
\begin{enumerate}\setlength\itemsep{3pt}
\item \textbf{Agents}: Who makes decisions? (Individuals, firms, government)
\item \textbf{Objectives}: What do agents care about? (Utility, profits, social welfare)
\item \textbf{Constraints}: What limits their choices? (Budget, technology, information)
\item \textbf{Environment}: What are the rules? (Tax system, market structure, policies)
\item \textbf{Equilibrium}: What happens when everyone optimizes?
\end{enumerate} 
\end{frame}



\begin{frame}
\frametitle{Constrained optimization: The workhorse of economic modeling}
Many economic problems can be written as constrained optimization:
\begin{equation*}
\max_{x_1, x_2, \ldots, x_n} f(x_1, x_2, \ldots, x_n)
\end{equation*}
\begin{equation*}
\text{subject to (s.t.): } g(x_1, x_2, \ldots, x_n) \leq c
\end{equation*}\vspace{-15pt}
\begin{itemize}\setlength\itemsep{3pt}
\item \textbf{Objective function} $f(\cdot)$: What the agent wants to maximize
\item \textbf{Choice variables} $x_1, \ldots, x_n$: What the agent can control
\item \textbf{Constraint} $g(\cdot) \leq c$: Limits on feasible choices
\end{itemize} \medskip

How calculus helps us solve this:
\begin{itemize}\setlength\itemsep{3pt}
\item Take derivatives: $\frac{\partial f}{\partial x_i}$ tells us how $f$ changes with small change in $x_i$
\item At optimum: Marginal benefit of increasing $x_i$ = Marginal cost (from constraint)
\item First-order conditions (FOCs): Set derivatives equal to zero (or use Lagrangian)
\item Solve system of equations to find optimal $x_1^*, \ldots, x_n^*$
\end{itemize}
\end{frame}

\begin{frame}
\frametitle{The Lagrangian method for constrained optimization}
The problem:
\begin{equation*}
\max_{x_1, \ldots, x_n} f(x_1, \ldots, x_n) \quad \text{subject to: } g(x_1, \ldots, x_n) = c
\end{equation*}

The Lagrangian approach: Convert constrained problem into unconstrained one
\begin{equation*}
\mathcal{L}(x_1, \ldots, x_n, \lambda) = f(x_1, \ldots, x_n) + \lambda \left[c-g(x_1, \ldots, x_n)\right]
\end{equation*}
where $\lambda$ is the \textbf{Lagrange multiplier}  \medskip

Solution method for $x_i^*$ and $\lambda^*$: Take derivatives and set equal to zero
\begin{itemize}\setlength\itemsep{3pt}
\item FOC for $x_i$: $\frac{\partial \mathcal{L}}{\partial x_i} = \frac{\partial f}{\partial x_i} - \lambda \frac{\partial g}{\partial x_i} = 0$
\item FOC for $\lambda$: $\frac{\partial \mathcal{L}}{\partial \lambda} = g(x_1, \ldots, x_n) - c = 0$ (ensures constraint holds)
\item Solve system of $n+1$ equations for $x_1^*, \ldots, x_n^*, \lambda^*$
\end{itemize} \medskip

Interpretation: $\lambda^*$ = how much objective improves if constraint relaxed by one unit
\end{frame}

\begin{frame}
\frametitle{Example: Optimal tax to maximize social welfare}

Objective: Maximize social welfare function $W(\cdot)$
\begin{equation*}
\max_{t} W(U_1(t), U_2(t), \ldots, U_n(t))
\end{equation*}
where $U_i(t)$ is utility of individual $i$ given tax rate $t$ \medskip

Constraint: Government budget constraint (revenue needed for spending $G$)
\begin{equation*}
G \leq  \sum_{i=1}^n t \cdot w_i \cdot \ell_i(t) 
\end{equation*}

Trade-off
\begin{itemize}\setlength\itemsep{3pt}
\item Higher $t$ allows more redistribution (equity)
\item But reduces labor supply $\ell_i(t)$ and distorts economy (efficiency)
\item Optimal tax balances these concerns
\end{itemize}
\end{frame}

\begin{frame}
\frametitle{From models to empirics: Why we need both}
Models provide structure to empirical thought
\begin{itemize}\setlength\itemsep{3pt}
\item What parameters matter? (e.g., labor supply elasticity, social preferences)
\item What would we observe if the model is correct?
\item How to extrapolate from one setting to another
\end{itemize} \medskip

Data are used to test and discipline models
\begin{itemize}\setlength\itemsep{3pt}
\item Do people actually respond to taxes as predicted?
\item How large are the responses? (quantitative magnitudes)
\item Are there important factors the model missed?
\end{itemize} \medskip


This course emphasizes \textit{both} theoretical frameworks and empirical evidence
\end{frame}

\begin{frame}
\frametitle{We will also use a lot of data and softwares} 
 Working with numerical data involves working with statistical softwares
\begin{itemize}\setlength\itemsep{3pt}
\item Process raw data: Excel, GIS programs (maps to numbers), OCR (digitalize text data)
\item Configure the data to a clean estimation-friendly form
\item Then run estimation to obtain results and make rigourous claims
\item The last two are done with Python, R, Stata ...
\end{itemize} \medskip
This involves many skill sets
\begin{itemize}\setlength\itemsep{3pt}
\item Understanding statistical inference
\item Interpreting intuitions from the numerical results
\item Patience (this is about 70-80\%)
\end{itemize} \medskip
\end{frame}

\begin{frame}
\frametitle{Next week} 
 We will cover classical public economics topics from week 3
\begin{itemize}\setlength\itemsep{3pt}
\item We plan to cover roughly 6 papers each course
\item I will assume that you have read at least the abstract and introduction of each papers
\begin{itemize}
  \item Ideally, try to read from cover to cover for 1-2 papers that interest you the most
\end{itemize}
\end{itemize} \medskip
What we will do: Econometrics 
\begin{itemize}\setlength\itemsep{3pt}
\item Cover applied econometrics theory and technics required to understand the papers
\item We will go over regression methods: OLS, Instrumental Variables, Panel regression, Difference-in-Differences, Regression discontinuity
\item Then, we will try demonstrations (demo files to be uploaded)
\begin{itemize}
  \item  If you have a laptop with R or Stata, I highly recommend bringing it to class
\end{itemize}
\end{itemize} \medskip
\end{frame}


%%%%%%%%%%%
\end{document}
